% Supplementary material for the journal paper (compile on its own with: latexmk -g -xelatex supplementary.tex).
% It reuses the preamble of main.tex and reads the labels of main.tex through xr-hyper (main.aux must exist: compile main.tex first).
\documentclass[preprint,12pt]{elsarticle}
% Shared preamble of main.tex and supplementary.tex (everything between \documentclass and \begin{document}).
\usepackage[english,bidi=default]{babel}
\babelfont{rm}{TeXGyreTermesX}
\babelprovide[import]{persian}
\babelfont[persian]{rm}[Scale=0.92]{Vazirmatn}
\newcommand{\fa}[1]{\foreignlanguage{persian}{#1}}
\usepackage{graphicx}
\usepackage{booktabs}
\usepackage{amsmath}
\usepackage{amssymb}
\usepackage{array}
\usepackage{multirow}
\usepackage{xcolor}
\usepackage{colortbl}
\usepackage{tabularx}
\usepackage{placeins}
\usepackage{enumitem}
\setlist{nosep,leftmargin=*}
\usepackage{pgfplots}
\pgfplotsset{compat=1.17}
\usepackage[most]{tcolorbox}
% Shared figure style for every figure in the paper (input once in the preamble of main.tex).
% One colour and one marker per pipeline, used everywhere (colour-blind-safe; markers keep figures legible in grey):
%   Hybrid grey / circle, HR teal / square, Lex-RAG orange / triangle, SetR-p+fill blue / diamond.
% Budget encoding where both budgets appear: open marker = 6,000 characters, filled marker = 15,500 characters.
\usetikzlibrary{positioning,arrows.meta,fit,backgrounds,calc,shapes.geometric,shapes.misc,shadows}
\usepgfplotslibrary{groupplots}
\definecolor{cHyb}{HTML}{6E6E6E}
\definecolor{cHR}{HTML}{009E73}
\definecolor{cLex}{HTML}{E66100}
\definecolor{cSetR}{HTML}{0072B2}
\definecolor{cInk}{HTML}{222222}
\definecolor{cAxis}{HTML}{5A5A5A}
\definecolor{cGrid}{HTML}{E4E4E4}
\definecolor{cHigh}{HTML}{FFF1E0}% light orange band behind the key comparison
\definecolor{cWithdrawn}{HTML}{B8B8B8}
\colorlet{cHybL}{cHyb!14}
\colorlet{cHRL}{cHR!14}
\colorlet{cLexL}{cLex!14}
\colorlet{cSetRL}{cSetR!14}
% legacy names (older figure sources) map onto the same palette
\colorlet{cbBlack}{cHyb}\colorlet{cbBlue}{cHR}\colorlet{cbVerm}{cLex}\colorlet{cbGreen}{cSetR}

\pgfplotsset{
  % axis look shared by all plots: light grid, thin grey axes, readable fonts (>= 8pt at print size)
  paper/.style={
    axis line style={cAxis, line width=0.6pt},
    tick style={cAxis, line width=0.6pt},
    major tick length=3pt,
    tick label style={font=\footnotesize, text=cInk},
    label style={font=\small, text=cInk},
    title style={font=\small\bfseries, text=cInk},
    grid style={cGrid, line width=0.5pt},
    legend style={draw=none, fill=none, font=\footnotesize},
    every axis plot/.append style={line width=1.5pt, mark size=3.2pt},
    axis x line*=bottom, axis y line*=left,
  },
  % pipeline identities: filled marker (single-budget plots and 15,500) and open marker (6,000)
  pHyb/.style={color=cHyb, mark=*, mark options={solid, fill=cHyb, draw=white, line width=0.5pt}},
  pHR/.style={color=cHR, mark=square*, mark options={solid, fill=cHR, draw=white, line width=0.5pt}},
  pLex/.style={color=cLex, mark=triangle*, mark options={solid, fill=cLex, draw=white, line width=0.5pt}, mark size=3.8pt},
  pSetR/.style={color=cSetR, mark=diamond*, mark options={solid, fill=cSetR, draw=white, line width=0.5pt}, mark size=3.8pt},
  pHybO/.style={color=cHyb, mark=*, mark options={solid, fill=white, draw=cHyb, line width=1.3pt}},
  pHRO/.style={color=cHR, mark=square*, mark options={solid, fill=white, draw=cHR, line width=1.3pt}},
  pLexO/.style={color=cLex, mark=triangle*, mark options={solid, fill=white, draw=cLex, line width=1.3pt}, mark size=3.8pt},
  pSetRO/.style={color=cSetR, mark=diamond*, mark options={solid, fill=white, draw=cSetR, line width=1.3pt}, mark size=3.8pt},
  % select rows of a data file by its numeric class column (1 Hybrid, 2 HR, 3 Lex-RAG, 4 SetR-p+fill)
  cls/.style={restrict expr to domain={\thisrow{cls}}{#1:#1}},
  % interval encodings: thick bar = identification bounds over unresolved judgments; thin whisker = sampling envelope
  boundbar/.style={only marks, mark=none, forget plot, error bars/.cd, x dir=both, x explicit,
    error bar style={line width=3.4pt, line cap=round}, error mark=none},
  ciwhisker/.style={only marks, mark=none, forget plot, error bars/.cd, x dir=both, x explicit,
    error bar style={line width=0.9pt}, error mark=|, error mark options={line width=0.9pt, mark size=2.6pt}},
  ybound/.style={only marks, mark=none, forget plot, error bars/.cd, y dir=both, y explicit,
    error bar style={line width=2.6pt, line cap=round}, error mark=none},
  xenvelope/.style={only marks, mark=none, forget plot, error bars/.cd, x dir=both, x explicit,
    error bar style={line width=0.9pt}, error mark=|, error mark options={line width=0.9pt, mark size=2pt}},
}
\pgfplotstableset{col sep=semicolon}

\tikzset{
  % direct labels next to data
  dlabel/.style={font=\footnotesize\bfseries, inner sep=1pt},
  note/.style={font=\footnotesize, text=cInk, align=left, inner sep=1pt},
  % pipeline diagram
  stage/.style={draw=#1, line width=0.9pt, rounded corners=4pt, fill=#1!10, align=center, font=\footnotesize,
    inner sep=3pt, minimum height=0.95cm, text=cInk},
  stage/.default=cHyb,
  shared/.style={draw=cAxis, line width=0.8pt, rounded corners=4pt, fill=white, align=center, font=\footnotesize,
    inner sep=3pt, minimum height=0.95cm, text=cInk, drop shadow={opacity=0.22, shadow xshift=0.7pt, shadow yshift=-0.7pt}},
  flow/.style={-{Stealth[length=5pt, width=4pt]}, line width=1.1pt, draw=#1},
  flow/.default=cAxis,
  fallback/.style={-{Stealth[length=4pt]}, line width=0.9pt, draw=#1, densely dashed},
  fallback/.default=cAxis,
  lanelabel/.style={font=\small\bfseries, text=#1, anchor=east, align=right},
  badge/.style={circle, fill=#1, text=white, font=\scriptsize\bfseries, inner sep=0pt, minimum size=11pt},
}
% class-indexed pipeline styles (for \pgfplotsinvokeforeach over the cls column)
\pgfplotsset{
  pcls1/.style={pHyb}, pcls2/.style={pHR}, pcls3/.style={pLex}, pcls4/.style={pSetR},
  pclsO1/.style={pHybO}, pclsO2/.style={pHRO}, pclsO3/.style={pLexO}, pclsO4/.style={pSetRO},
}
\colorlet{pcol1}{cHyb}\colorlet{pcol2}{cHR}\colorlet{pcol3}{cLex}\colorlet{pcol4}{cSetR}
% inline pipeline key: \pkey{colour}{mark}{name} draws the pipeline's marker followed by its name in its colour
\newcommand{\pkey}[3]{\tikz[baseline=-0.55ex]{\pgfsetplotmarksize{3.2pt}\pgfsetfillcolor{#1}\pgfsetstrokecolor{#1}\pgfsetlinewidth{0.5pt}\pgfuseplotmark{#2}}\,\textcolor{#1}{\textbf{#3}}}
\newcommand{\pkeyopen}[3]{\tikz[baseline=-0.55ex]{\pgfsetplotmarksize{3.2pt}\pgfsetfillcolor{white}\pgfsetstrokecolor{#1}\pgfsetlinewidth{1.3pt}\pgfuseplotmark{#2}}\,\textcolor{#1}{#3}}
\newcommand{\pipelinekey}{\pkey{cHyb}{*}{\hybrid}\quad\pkey{cHR}{square*}{\hr}\quad\pkey{cLex}{triangle*}{\lexb}\quad\pkey{cSetR}{diamond*}{\setr}}
\newcommand{\budgetkey}{\pkeyopen{cAxis}{*}{6,000 chars}\quad\pkey{cAxis}{*}{\mdseries 15,500 chars}}

% post-hoc additions (rewrite 2026-11): Lex-RAG without selection (reddish purple / diamond-open), search agents (vermillion),
% evidence-reach shades (sky blue = answer correct, vermillion = answer wrong).
\definecolor{cFus}{HTML}{CC79A7}
\definecolor{cAgent}{HTML}{D55E00}
\definecolor{cHit}{HTML}{56B4E9}
\definecolor{cMiss}{HTML}{D55E00}
\pgfplotsset{
  pFus/.style={color=cFus, mark=pentagon*, mark options={solid, fill=cFus, draw=white, line width=0.5pt}, mark size=3.6pt, densely dashed},
  pcls5/.style={pFus}, pclsO5/.style={color=cFus, mark=pentagon*, mark options={solid, fill=white, draw=cFus, line width=1.3pt}, mark size=3.6pt},
  pcls6/.style={color=cAgent, mark=square*, mark options={solid, fill=cAgent, draw=white}}, pclsO6/.style={pcls6},
  pcls7/.style={color=cAxis, mark=triangle*, mark options={solid, fill=cAxis, draw=white}}, pclsO7/.style={pcls7},
}
\colorlet{pcol5}{cFus}\colorlet{pcol6}{cAgent}\colorlet{pcol7}{cAxis}

\usepackage{listings}
\lstset{basicstyle=\ttfamily\scriptsize, breaklines=true, columns=fullflexible, frame=single, framesep=2pt}
% \usepackage{lineno} % optional
\usepackage[hidelinks]{hyperref}
\biboptions{numbers,sort&compress}
\setlength{\emergencystretch}{2em}
\renewcommand{\topfraction}{0.9}
\renewcommand{\bottomfraction}{0.9}
\renewcommand{\textfraction}{0.1}
\renewcommand{\floatpagefraction}{0.85}
\setcounter{topnumber}{3}
\setcounter{bottomnumber}{3}
\setcounter{totalnumber}{5}
\renewcommand{\figurename}{Fig.}
\AtBeginDocument{\renewcommand{\figurename}{Fig.}}% elsarticle and babel reset it at the start of the document
\addto\captionsenglish{\renewcommand{\figurename}{Fig.}}
\setlength{\belowcaptionskip}{4pt}
\newcolumntype{L}[1]{>{\raggedright\arraybackslash}p{#1}}
% system names (same macros as the earlier version)
\newcommand{\hybrid}{\textsc{Hybrid}}
\newcommand{\hr}{\textsc{HR}}
\newcommand{\lexb}{\textsc{Lex-RAG}}
\newcommand{\setr}{\textsc{SetR-p+fill}}
\newcommand{\ftwo}{three-signal score fusion}

\usepackage{xr-hyper}
\externaldocument{main}
\addtolength{\textheight}{14pt}% the long log and settings tables are floats of nearly a full page
\renewcommand{\thesection}{S\arabic{section}}
\renewcommand{\thesubsection}{S\arabic{section}.\arabic{subsection}}
\renewcommand{\thetable}{S\arabic{table}}
\renewcommand{\thefigure}{S\arabic{figure}}

\begin{document}
\begin{frontmatter}
\title{Supplementary material for: What Reaches the Reader: A Comparison of Retrieval-Augmented Generation Setups in Persian and English}
\author[ut]{Arash Nasr Esfahani\corref{cor1}}
\ead{arash.nasr@ut.ac.ir}
\cortext[cor1]{Corresponding author.}
\affiliation[ut]{organization={School of Mathematics, Statistics, and Computer Science, University of Tehran},city={Tehran},country={Iran}}
\begin{abstract}
This document holds the post hoc log and settings, the development record, the prompts of the paper and the examination of the gold-evidence misses. Section numbers start with S, and references to Sections, Tables and Appendices without S point to the main paper.
\end{abstract}
\end{frontmatter}

\section{Post hoc variants: log and settings}
\label{supp:log}
% src: moved from Appendix E of the main paper (Tables E.1 to E.4 and the run log) on 2026-11-18; work/posthoc_2026-11/COMMON.md and S1..S15 RESULTS.md
Table~\ref{tab:posthoc-log} lists every variant and analysis that was run after the frozen test, including those that the body of the paper does not discuss.
Costs are as reported by the provider, and a dash marks costs that were not itemised.
Tables~\ref{tab:posthoc-settings}, \ref{tab:posthoc-settings2} and~\ref{tab:posthoc-settings3} give the settings of each variant.
The rules that held for every run are stated in \ref{app:posthoc} of the main paper.
Two runs were cut short and are not used.
One was a first pass of the first judge that was stopped after five verdicts in order to run the variants in parallel.
The other was a first selector-swap attempt that failed on a software error in the runner and was repeated from scratch.
Review packets for a second reviewer were prepared, but no human verdicts exist for them.
The second judge scored seven of the variants: \lexb{} without selection, the selected set of \setr{} without fill, the listwise reranker, the BM25 agent, the regular-expression agent, gold evidence and the closed-book run.
Its scores are 57 to 59, 63 to 64, 64, 66 to 68, 67 to 68, 82 to 84 and 26 to 28, against 58 to 59, 63, 66 to 67, 73, 68, 83 and 26 under the first judge.
% src: every row from work/posthoc_2026-11/S*/RESULTS.md (section 3 cost, section 4 failed or skipped, section 2 results); dates from the ledger.jsonl
%      timestamps of each stream (all 2026-11-10; the date column was dropped, REVIEW2 item 4.4); costs are sums of usage.cost in each stream's ledger (S1 0.8325; S2a 0.842 = reader 0.695 + judge 0.147;
%      S2b 1.3239; S3 0.274; S4 1.084; S6 0.1760; S5 and S7 no API calls). Per-arm costs where the RESULTS file gives them (S3: GREP test 0.1173, ITER test 0.1151;
%      S4: LC-GOLD 0.414, LC-BM25 0.435, LC-POOL 0.209, router 0.008; S6: CLOSEDBOOK reader 0.0085).
% Added 2026-11-11: S9 rows (S9_gpu/RESULTS.md sec. 3: cost 0.127, 77/90, 9,070 tokens; job 1 alignment: 16/96 HYBRID_6000 packs would change) and the S8 second-judge row (S8_factorial/tables.md); S10-S15 rows: S10 0.2763, S11 0.112, S12 0.1798, S13 0.0920, S15 0.078 (RESULTS.md sec. 3 of each); S14 and the spelling variants of S15 make no API call.
\begin{table*}[tp]
\caption{Post hoc. Every variant and analysis run after the frozen test}
\label{tab:posthoc-log}
\centering
\scriptsize
\setlength{\tabcolsep}{4pt}
\begin{tabularx}{\textwidth}{@{}p{4.4cm}rXr@{}}
\toprule
Variant or analysis & USD & Outcome & Reported in \\
\midrule
\multicolumn{4}{@{}l}{\emph{Packs and budgets (block cost 0.84)}} \\
Rebuild of the 576 frozen packs & 0 & All 576 match the saved packs byte for byte & \ref{app:posthoc} \\
\lexb{} without selection, five budgets & -- & 58--59 at 6,000 against 68--69 with selection; 65 against 67 at 15,500 & Sec.~\ref{sec:lost} \\
\hybrid{} and \hr{} at 3,000, 10,000, 24,000 & -- & \hybrid{} rises from 49--51 to 72; \hr{} is flat from 6,000 and falls to 66 at 24,000 & Sec.~\ref{sec:budget} \\
Gold evidence only & -- & 83 of 90 & Sec.~\ref{sec:lost} \\
Examination of the seven gold-evidence misses & 0 & Four bundles that do not tie the key value to its subject, two answers failed by a judge, one incomplete answer & Sect.~\ref{supp:goldmiss} \\
MMR packer & -- & 60; does not replace the selector & Sec.~\ref{sec:lost} \\
Knapsack packer & -- & 44; far below the fusion order & Sec.~\ref{sec:lost} \\
\midrule
\multicolumn{4}{@{}l}{\emph{Other baselines (block cost 0.18)}} \\
Closed book & 0.01 & 26 of 90; 3 of 6 abstentions & Sec.~\ref{sec:lost} \\
Listwise LLM reranker & 0.11 & 66--67 (64 under the second judge); no difference from \hybrid{} or \hr{} & Sec.~\ref{sec:lost} \\
SetR, selected set only & 0.02 & 63, the same as \setr{} with fill & Sec.~\ref{sec:lost} \\
First judge pass of these variants & -- & Stopped after 5 verdicts to run in parallel; not used & -- \\
\midrule
\multicolumn{4}{@{}l}{\emph{Selection across budgets and candidate sources (block cost 0.28 for the first four steps)}} \\
\lexb{} at 3,000, 10,000, 24,000 (saved selection re-packed) & -- & 59, 69--70, 72; the gain over Lex-RAG without selection is large at 3,000 (+15..+16) and gone from 15,500 & Sec.~\ref{sec:budget} \\
Whole pipeline re-run with no selection deadline (all 96 questions) & -- & 71--72 against 68--69; the gain comes from the 19 saved fallbacks & Sec.~\ref{sec:lost} \\
Selection over the reranker's candidates & -- & 70; no gain over the reranker's own pack & Sec.~\ref{sec:lost} \\
\lexb{} without selection on the 60 development questions & -- & 29 against 37 for \lexb{} on the 54 answerable ones (+8, sign $p$ .021) & \ref{app:posthoc} \\
\midrule
\multicolumn{4}{@{}l}{\emph{Searching harder (block cost 1.36)}} \\
BM25 search agent & 0.12 & 73 of 90 at 9,285 tokens per question (66--68 under the second judge) & Sec.~\ref{sec:search} \\
Regular-expression agent & 0.12 & 68 of 90 at 12,338 tokens per question (67--68 under the second judge) & Sec.~\ref{sec:search} \\
Long context, gold documents & 0.41 & 79 of 90; uses the annotation & Sec.~\ref{sec:search} \\
Long context, BM25-ranked documents & 0.44 & 24 of 90; the cut text lost most of the evidence & Sec.~\ref{sec:search} \\
Long context, whole pool & 0.21 & Three questions only (cost); all correct; only two of six pools fit the context & \ref{app:posthoc} \\
Routing by tier, by a router model, by answer format & 0.01 & Cheaper by about a fifth at best, and only with the annotated tier & Sec.~\ref{sec:search} \\
\bottomrule
\end{tabularx}
\end{table*}

\begin{table*}[tp]
\addtocounter{table}{-1}
\renewcommand{\theHtable}{\thetable.cont}
\caption{Post hoc. Every variant and analysis run after the frozen test (continued)}
\centering
\scriptsize
\setlength{\tabcolsep}{4pt}
\begin{tabularx}{\textwidth}{@{}p{4.4cm}rXr@{}}
\toprule
Variant or analysis & USD & Outcome & Reported in \\
\midrule
\multicolumn{4}{@{}l}{\emph{Hybrid-index agent and vector check (block cost 0.13)}} \\
Hybrid-index search agent (reranker as search tool) & 0.13 & 77 of 90 at 9,070 tokens per question; 6 to 8 above \hr{}, not significant & Sec.~\ref{sec:search} \\
Ollama against Hugging Face bge-m3 vectors & 0 & The same model up to numerical noise; with the other vectors 16 of 96 \hybrid{} packs at 6,000 characters would change & \ref{app:posthoc} \\
Second judge on the post hoc variants (Mistral Large) & 1.22 & 19 variants scored: seven first (the BM25 agent 66--68 against 73 under the first judge, the listwise reranker 64 against 63 for \hr{}), then the hybrid-index agent (74--75), the ten component variants at 6,000 characters and the cross-language run (92\% agreement, $\kappa=0.82$); the budget sweep, routing, long context and the selector-over-reranker and no-deadline variants were not scored & Sec.~\ref{sec:stability} \ % src: S16_judge2/tables.md, ledger_j2.jsonl (1.2248 USD)
\midrule
\multicolumn{4}{@{}l}{\emph{Retrieval components (block cost 0.66)}} \\
Dense only and BM25 only, 6,000 and 15,500 characters & 0.28 & 53 and 48 at 6,000, 63 and 61 at 15,500, against \hybrid{} 62--63 and 68; no GPU used & Sec.~\ref{sec:components} \\
HyDE passage and multi-query rewriting & 0.11 & 61 and 58--59 against 62--63; fewer complete bundles in Persian & Sec.~\ref{sec:components} \\
Chunks of 500 and 2,000 characters, \hybrid{} and \hr{} & 0.18 & 1,000 best; 2,000 loses 13--14 (\hybrid{}) and 18--20 (\hr{}); two indexes of 119,308 and 28,205 chunks embedded & Sec.~\ref{sec:components} \\
multilingual-e5-large in place of bge-m3, hybrid and dense only & 0.09 & 58 and 54--55 against 62--63 and 53; 58,218 chunks embedded in 17.6 minutes of GPU time & Sec.~\ref{sec:components} \\
\midrule
\multicolumn{4}{@{}l}{\emph{Persian and English (block cost 0.08)}} \\
Questions asked in the other language (translation, four retrieval configurations, \hr{}, closed book) & 0.08 & BM25 finds a complete bundle for 2 and 8 of 45; \hr{} 53--54 against 69--71 & Sec.~\ref{sec:language} \\
Spelling variants of the 45 Persian questions (retrieval only) & 0 & Normalised retrieval loses at most 1 bundle; BM25 without normalisation loses 5 to 8 & Sec.~\ref{sec:language} \\
Language splits, permutation tests, tokens per reader, judges by language (saved data) & 0 & Regular-expression agent split survives correction; tokens per character about 1.5 to 3.0 times higher in Persian & Sec.~\ref{sec:language} \\
\midrule
\multicolumn{4}{@{}l}{\emph{Stability of the ranking (block cost 2.16: second judge 0.83, variance runs 1.32)}} \\
Second judge (Mistral Large) on 768 answers & 0.83 & 91.2\% agreement; the last place of \setr{} does not reproduce & Sec.~\ref{sec:stability} \\
Human-review packets for a second reviewer & 0 & Built and tested on synthetic data; no human verdicts yet & \ref{app:posthoc} \\
Reader repeated twice on identical packs (1,536 answers) & -- & Standard deviation about two answers per configuration & Sec.~\ref{sec:stability} \\
Selector re-sent twice (\lexb{}) and once (\setr{}) & -- & The pack changes for 44--50 (\lexb{}) and 52 (\setr{}) of 96 questions; accuracy within reader noise & Sec.~\ref{sec:stability} \\
Selector swapped for a smaller model & -- & 59; 40 of 96 selections fall back, 21 of them for invalid passage identifiers & Sec.~\ref{sec:lost} \\
Selection call re-run without the deadline (the 15 timed-out calls) & -- & 14 of 15 valid replies; 7 against 5 correct & \ref{app:variance} \\
One selection packed at both budgets & -- & No budget gain for \lexb{} either way & \ref{app:variance} \\
Discarded first run of the selector swap & 0.003 & Harness error; rows discarded, run repeated from scratch & -- \\
\midrule
\multicolumn{4}{@{}l}{\emph{Analyses of saved data (no model calls)}} \\
Mixed models, equivalence tests, detectable effect & 0 & Pooled effect of selection over \hybrid{}; no reader interaction & Sec.~\ref{sec:stability} \\
Earlier non-adaptive mixed model & 0 & Did not converge in the variance terms; superseded & -- \\
Evidence reach, contamination, tokens, errors, Persian channels & 0 & See \ref{app:analysis} & Sec.~\ref{sec:lost} \\
\bottomrule
\end{tabularx}
\end{table*}

% src: work/posthoc_2026-11/S2a_packs/RESULTS.md sec. 1 (FUSION, GOLD, MMR, KNAP, budgets); S6_baselines/RESULTS.md sec. 1 (CLOSEDBOOK, LISTWISE, SETR_NOFILL);
%      S3_agentic/RESULTS.md sec. 1 (agents); S4_route_long/RESULTS.md sec. 1 (routing, long context); S2b_variance/RESULTS.md sec. 1 (variance runs);
%      S1_evaluator/RESULTS.md sec. 1 (second judge). Reader GLM-5.3-flash (Together pin) and judge gpt-oss-120b (deepinfra/bf16, temperature 0, seed 42) as in the frozen test.
\begin{table*}[tp]
\caption{Post hoc. Settings of the pack, baseline and judge variants}
\label{tab:posthoc-settings}
\centering
\scriptsize
\setlength{\tabcolsep}{3pt}
\begin{tabularx}{\textwidth}{@{}p{2.6cm}p{2.9cm}Xp{2.8cm}@{}}
\toprule
Variant & Input & Procedure & Settings \\
\midrule
\lexb{} without selection & The 40 fused evidence units of \lexb{} & Whole units in fusion order until the exact rendered cost reaches the budget; never clipped; no LLM call & Budgets 3,000; 6,000; 10,000; 15,500; 24,000 characters \\
\hybrid{}, \hr{} at new budgets & Saved top 40 (fused; reranked) & The frozen packing rule at each budget & 3,000; 10,000; 24,000 characters \\
Gold evidence & First annotated evidence bundle & Quote text taken from the corpus offsets, one item per evidence unit, in bundle order & Mean 572 characters; 90 answerable questions \\
MMR diversity rule & The same 40 units & Maximal marginal relevance: relevance is the min-max normalised fusion score, redundancy the largest cosine to units already chosen (saved bge-m3 embeddings); whole units only & $\lambda=0.7$, set a priori; 6,000 characters \\
Knapsack rule & The same 40 units & Greedy by fusion score per rendered character; whole units only & 6,000 characters \\
Closed book & The question only & The answer prompt without evidence: ``answer from your own knowledge'', otherwise \texttt{NOT\_FOUND} & One call \\
Listwise LLM reranker & \hybrid{}'s top 20 chunks & One RankGPT call with the published multi-turn prompt; passages cut to 700 characters; chunks 21--40 follow in \hybrid{} order; the packer then fills as for \hybrid{} & Reasoning effort low; 2,000 tokens; no deadline \\
SetR, selected set only & \setr{}'s saved selected set & The set is offered alone to the packer, without the \hybrid{} fill; where \setr{} had fallen back (6 questions), its own fallback pack & Mean 3.04 chunks, 3,056 characters (against 5,949 with fill) \\
Second judge & The 768 saved GLM answers & Same rubric, prompt, parser and retries as judge 1; a reply that does not parse after three tries is unresolved & Mistral Large (2512), provider \texttt{mistral/zdr}, temperature 0, seed 42 \\
\bottomrule
\end{tabularx}
\end{table*}

\begin{table*}[tp]
\caption{Post hoc. Settings of the search, long-context, routing and variance variants}
\label{tab:posthoc-settings2}
\centering
\scriptsize
\setlength{\tabcolsep}{3pt}
\begin{tabularx}{\textwidth}{@{}p{2.6cm}p{2.9cm}Xp{2.8cm}@{}}
\toprule
Variant & Input & Procedure & Settings \\
\midrule
BM25 search agent & The frozen 1,000-character chunks of the question's pool & Tool \texttt{search(query)} returns the top 8 chunks (BM25 channel only, CPU); \texttt{finish(chunk ids)} hands evidence to the answer call, packed at 15,500 characters & GLM-5.3-flash; at most 5 rounds, 3 tool calls per round; median evidence 2,735 characters \\
Regular-expression agent & Raw document text of the pool; documents named D1, D2, \ldots & Tools \texttt{grep} (regular expression, at most 20 lines), \texttt{read} (at most 2,500 characters), \texttt{list\_docs}, \texttt{finish} (spans of at most 4,000 characters); Persian letters and digits folded to one form & As above; median evidence 1,652 characters \\
Long context, gold documents & The documents that hold the gold evidence & Equal share of the prompt per document, cut to a window around the gold units, so every gold unit is shown (228 of 228) & 32,000 prompt tokens; 90 questions \\
Long context, BM25 documents & Pool documents ranked by document-level BM25 & Whole documents added until the limit; the first that does not fit is cut to its head & 32,000 prompt tokens; 96 questions; gold units kept: 62 of 228 \\
Long context, whole pool & The whole pool & Only the two LAW pools fit the 1,048,575-token context & 3 questions (LAW-fa), 464,389 to 464,455 prompt tokens \\
Router model & The question text only & Reply A (one fact, route to \hybrid{} 6,000) or B (combine evidence, route to the wide pipeline); the prompt was chosen on dev out of three & 600 tokens; 193 tokens per question \\
Router by tier, by answer format & Annotated tier; the answer format already given to every pipeline & Tier 1 to \hybrid{} at 6,000, tiers 2 and 3 to the wide pipeline; the format rule sends single-fact answers to tier 1 & No model call \\
Reader repeated & The identical saved pack messages of all 8 configurations & Sent again twice to the same reader and judged & 1,536 answers \\
Selection call re-sent & The saved selector prompt & Re-sent to the GLM selector with the frozen settings, re-packed; unchanged packs reuse the matching reader answer & \lexb{} twice, \setr{} once \\
Selector swapped & The same selector prompt & \texttt{gemini-3.1-flash-lite} chooses the passages; the GLM reader answers & Same deadline, cap and strict grammar \\
Selection call without deadline (15 timed-out calls) & The 15 questions whose selector call timed out & The GLM selector with no deadline, then re-pack and answer & 600-token cap kept \\
One selection, both budgets & The saved 6,000 and 15,500 selector replies & Each reply packed at the other budget & 82 answers \\
\lexb{} at other budgets & The saved 6,000-character selector reply & Re-packed with the frozen code at the new budget; the selector prompt does not depend on the budget (checked byte for byte), so there are no new selection calls & 3,000; 10,000; 24,000 characters \\
Whole pipeline without deadline (all 96 questions) & The saved selector prompt, all 96 questions & Re-sent to the GLM selector with no deadline, re-packed and answered; an unchanged pack reuses the frozen answer & 600-token cap; 14 of 96 calls took over 3 seconds \\
Selection over the reranker's candidates & The reranker's top 24 in the frozen selector view & The same selection call without deadline, strict grammar, exact bundle packing, and fill and fallback in reranked order & 6,000 characters; 4 fallbacks \\
\lexb{} without selection on development data & The 60 development questions, saved \lexb{} traces & The pack without selection rebuilt from the saved candidates (equals the saved fallback pack in 60 of 60), answered and judged as on test & 54 answerable questions \\
\bottomrule
\end{tabularx}
\end{table*}

% src: work/posthoc_2026-11/S10_retrieval/RESULTS.md sec. 1; S11_rewrite sec. 1; S12_chunks sec. 1; S13_embed sec. 1; S14_bilingual sec. 1; S15_crossling sec. 1. Reader GLM-5.3-flash (Together pin), judge gpt-oss-120b as in the frozen test.
\begin{table*}[tp]
\caption{Post hoc. Settings of the retrieval-component and language variants}
\label{tab:posthoc-settings3}
\centering
\scriptsize
\setlength{\tabcolsep}{3pt}
\begin{tabularx}{\textwidth}{@{}p{2.6cm}p{2.9cm}Xp{2.8cm}@{}}
\toprule
Variant & Input & Procedure & Settings \\
\midrule
Dense only, BM25 only & The frozen \hybrid{} index (chunk vectors, BM25) and the saved query vectors & Top 40 of one channel, which keeps its order through rank fusion; \hybrid{}'s packer and answer prompt unchanged; the BM25 pack at 6,000 equals a frozen pack for one question and reuses its verdict & 6,000 and 15,500 characters; no new embedding, no GPU; 383 reader calls \\
HyDE & The question and the genre of its pool & One LLM call writes a hypothetical passage of 60 to 120 words in the question's language; dense search uses the passage vector, BM25 the question; rank fusion and the frozen packer follow & One passage, not averaged with the question vector; mean 649 characters; 6,000 characters \\
Multi-query & The question & One LLM call writes three paraphrases or sub-queries; BM25 and dense lists for the original and each rewrite (eight lists) are merged by rank fusion and the top 40 packed & All 96 questions received three rewrites; 6,000 characters \\
Chunk size & The documents of the six pools & New indexes with chunks of 500 (overlap 100) and 2,000 (overlap 400) characters; same BM25, bge-m3, fusion and reranker; the 1,000-character variants are the frozen cells, reproduced byte for byte & 119,308 and 28,205 chunks; packer top 8 and 40 candidates unchanged; 6,000 characters \\
Second embedding model & All 58,218 chunks and the questions & multilingual-e5-large with the model's passage and query prefixes, fp16, mean pooling; dense only, and fused with the frozen BM25; frozen packer & Commit 3d7cfbd; 512 tokens (65 chunks truncated); 6,000 characters \\
Translated questions & The 96 test questions & One plain translation call per question, fixed on five dev questions, with names and numbers kept (numbers identical for 96 of 96); BM25, dense, fused and reranked retrieval against the original pool; \hr{} and closed book answer in the language asked, and the judge sees the original question & 20 of 45 translations of English-pool questions keep Latin-script names; 6,000 characters \\
Spelling variants & The 45 Persian test questions & Arabic letters in place of Persian ones, zero-width non-joiner removed or replaced by a space, digits set to Persian, Arabic-Indic or Latin, and one combined variant (Arabic letters, non-joiner removed, Arabic-Indic digits); BM25 with the frozen and with a plain tokenizer, dense with and without normalisation, and fused & Retrieval only, no reader calls; top 40 and 6,000-character pack \\
Language splits and judges by language & Saved verdicts, ledgers and packs & Net answers and an exact sign test per half; permutation test of language by variant over the 20 source groups, with Holm correction over 14 contrasts; prompt tokens per evidence character by regression for six readers; judge agreement and kappa by language & No model calls; second judge as a snapshot of its run \\
\bottomrule
\end{tabularx}
\end{table*}


\section{Development record}
\label{app:negative}
% src: lex5_pilot/PILOT.md per-arm table (LEX5 8/12 9146 tok/q; LEX5+selector=glm 6/12 7496; HYBRID 2143); devrun/RESULTS_DEV.md (LLM calls/q 3.0); budget_matrix/BUDGET_MATRIX.md (HYBRID+KW 10, DOC-3 11)
% src (per row): devrun/RESULTS_DEV.md (LEX2, LEX3, AGENTIC; local gemma4:e4b reader); diagnosis/LEX3_PROFILE.md; lex4_replay/REPLAY.md; lex5_replay/REPLAY5.md; lex5_pilot/PILOT.md; lex5b/LEX5B.md; lex5b_fast/LEX5B_FAST.md;
%   lex6/eval/LEX6_EVAL.md; idea_test/IDEAS.md; lex_research/OFFLINE_TESTS.md; lex7/LEX7_OFFLINE.md; lex8/LEX8.md; lex9/LEX9.md; lex10/LEX10.md; lex11/LEX11.md; budget_matrix/BUDGET_MATRIX.md
% src (learned lasso term weights row): work/analysis/LEX_LEARNED.md (Result 1: recall@8 hand 0.738, 1-SE 0.726; CV fits dense, 230-284 nonzero of ~290; sparse regime: hand-set + sparse corrections 0.738); work/runs/lex_learn/compare/COMPARE.md
%   (wide dev acc n=63: hand-set 74.6, learned 68.3; paired 1 better / 5 worse / 57 equal); results.md there (gemma4:e4b reader, pilot.jsonl)
Table~\ref{tab:negative-full} lists every development variant that preceded or competed with \lexb{}.
All numbers in it are on the development set.
Retention (ret.) is decisive-span retention over the 54 answerable development questions at 6,000 characters, with the number of questions whose best bundle is fully retained in brackets (\hybrid{} 0.720 (32), \hr{} 0.806 (35)).
Screen is the number of strictly correct answers on a fixed 24-question subset with the GLM reader (\hybrid{} 12, \hr{} 13, \hybrid{} at 15,500 characters 19).
The first block used the earlier 60-question development run with a local 4B reader (gemma4:e4b), where \hybrid{} scored 27 and \hr{} 28 of 60.
Retention and screen numbers are not comparable with test accuracy.
Every admission rule was declared before the corresponding local result, but the sequence of variants was chosen in response to earlier results, so the whole record is exploratory.
The row marked with a dagger is an earlier experiment with an earlier version of the development set, a regex-based lexical method, a local reader and match-based accuracy on a wider set, and its retrieval is measured as gold-window recall.
The learned-weights experiment predates the benchmark.
All other variants were run between 25 September and 1 October 2026 by agents that executed the designs of the author (Table~\ref{tab:timeline}).
% src: moved from sections_old/F_development.tex
\begin{table*}[tp]
\caption{Every development variant with its main numbers}
\label{tab:negative-full}
\centering
\scriptsize
\begin{tabularx}{\textwidth}{@{}p{2.6cm}Xp{4.6cm}p{2.2cm}@{}}
\toprule
Variant & Idea & Result & Verdict \\
\midrule
Regex document scan & LLM plans keywords; a second call picks the documents' own phrases; code compiles Persian-aware regexes; structure windows & 22/60 (local reader); ret.\ 0.513 (20); 3.0 recorded LLM calls/q including the answer & below both baselines \\
Regex scan with routing & the regex scan + card routing, phrase mining, follow-up passes & 28/60 (= \hr); 3.0 calls/q; ret.\ 0.709 (30); p50 36.6\,s vs.\ 6.1\,s; 1 of 100 packed decisive spans found by the scan alone & 6$\times$ slower, no recall \\
\quad + cards/sections & document cards and section expansion & 24/60 & options hurt ($-4$) \\
Agentic loop (two settings) & plan, search, verify rows, answer (3.2--3.5 calls) & 25/60, 23/60; p50 43--49\,s & no gain, 7--8$\times$ slower \\
\midrule
Per-requirement retrieval (offline) & per-requirement retrieval, lexical votes, slot reservation, cover-then-fill & best (no lexical list) 0.718 (30); full design 0.681 (27); rule A 0.607 (22) & no variant passed \\
Listwise LLM selection (replay) & LLM listwise selection over \hybrid{} top-24 (full / snippet / head views) & full view 0.858 (42); snippet view 0.837 (39); head view 0.782 (33) & best retention so far \\
Full-passage selector (pilot, 12 q) & same, GLM reader; selector gpt-oss-20b or GLM & gpt-oss-20b selector 8/12, GLM selector 6/12, vs.\ \hybrid{} 5, \hr{} 7; tokens 4.3$\times$ and 3.5$\times$ \hybrid & failed 1.5$\times$ token gate \\
\lexb{} selector views & excerpt types under a 2,600-token view & smart 0.801, full top-10 0.792, literal 0.854 (38) & literal admitted \\
Faster selector & smaller view (top-16); gpt-oss-20b as selector & 0.813, 0.788; the gpt-oss-20b selector cut selector p50/p95 from 1,252/3,030 to 1,181/2,066\,ms but missed the required p50 cut and the retention gate & not admitted \\
\bottomrule
\end{tabularx}
\end{table*}

\begin{table*}[tp]
\addtocounter{table}{-1}
\renewcommand{\theHtable}{\thetable.cont}
\caption{Every development variant with its main numbers (continued)}
\centering
\scriptsize
\begin{tabularx}{\textwidth}{@{}p{2.6cm}Xp{4.6cm}p{2.2cm}@{}}
\toprule
Variant & Idea & Result & Verdict \\
\midrule
Keyword-count reranking; document-first routing (24 q) & keyword-count ranking; document-first search in top-$k$ docs & ret.\ (24 q) keyword-count only 0.601, \hybrid{} with keyword fusion 0.663, document-first (top-3) 0.642 vs.\ \hybrid{} 0.748; screen (answers): keyword fusion 10, document-first 11 & no \\
Document slot reservation & document ranking + slot reservation by evidence unit; cards; sections; router & 0.733 (32); sections 0.709; router: 54/54 deadline misses; screen 11 & stopped \\
bge-m3 channels (offline) & bge-m3 sparse / multi-vector channels, RRF, MMR, doc-decay & RRF of BM25, dense and sparse 0.763 (34), 4 better / 0 worse than \hybrid; multi-vector 0.730; MMR 0.767 & sparse useful, multi-vector not \\
Chunk prefixes; order fixer & contextual chunk prefixes; small LLM order fixer & prefixes 0.714 ($-0.049$); fixer 0.773 (35); fixer screen 13 & below \hr \\
Score-fusion variants & score fusion with learned sparse weights (CombSUM, MaxMean), compact units, coverage packer, synthetic-question field (\$3.49) & three-signal score fusion (MaxMean) 0.779 (33); CombSUM 0.746; coverage packer 0.712; +synthetic 0.726 & three-signal fusion kept as candidate generator \\
Learned ranker & LightGBM ranker trained on 18k synthetic questions & held-out nDCG@10 0.857 (vs.\ three-signal fusion 0.810); dev 0.761 (32) & no transfer \\
Wider pack without selection & spend token allowance on a wider pack without selection (8,500 chars) & 0.808 (36); not admitted (needs 38) & -- \\
\quad same at 15,500 & same at 15,500 chars & 0.882 (42); screen 15 vs.\ \hybrid@15,500 19 & below \hybrid@15,500 on the screen \\
Fusion membership, hybrid order & three-signal fusion membership, \hybrid{} presentation order & screen 12 @6,000, 15 @15,500 (= pack without selection) & this ordering intervention did not improve screen accuracy \\
\midrule
Learned lasso term weights$^\dagger$ & lasso over window features and phrase indicators replaces the hand-set window-scoring weights; five folds over questions, out of fold & accuracy 68.3\% vs.\ 74.6\% hand-set (63 paired records: 1 better / 5 worse / 57 equal); gold recall @8 windows 0.726 vs.\ 0.738; cross-validated fits dense (230--284 nonzero of about 290); sparse fits on top of the hand-set score 0.738 & not adopted; outside the theory's sparse-signal assumption \\
\bottomrule
\end{tabularx}

{\footnotesize\emph{Note.} $^\dagger$Earlier experiment on an earlier development set, described in the text.\par}
\end{table*}


Two patterns explain most rows.
Query-time lexical search failed because keywords written by a 20B model mostly paraphrased the question or invented wordings that the documents do not use.
Index-time enrichment, such as contextual prefixes and 141 thousand synthetic questions for \$3.49, did not move decisive chunks into the pack.
The one component that clearly raised retention was the learned-sparse channel.
Three-signal score fusion raised retention from 0.720 to 0.779 at 6,000 characters and from 0.847 to 0.882 at 15,500 characters. Yet its pack answered no more screen questions than \hybrid{} (12 of 24 at 6,000 characters, and 15 against 19 at 15,500).
Retention therefore does not explain accuracy here.
Cheaper selection did not replace the full selector.
In an offline replay of the replies of the admitted selector, keeping only the selected bundles removed approximately 60\% of the characters and lowered retention from 0.854 to 0.778, and keeping 20\% fill gave 0.815.
% src: lex11/LEX11.md selector pruning replay (S-c 5,766 chars 0.8540; core only 2,283 chars 0.7778; core+20% 3,230 chars 0.8148)
On the screen, \hybrid{} rose from 12 to 19 with the larger budget and \lexb{} fell from 17 to 15, but the selector hit its deadline in 10 of 24 questions at 15,500 characters (5 of 24 at 6,000). That budget effect is therefore confounded by provider latency.
The learned-weights experiment used lasso weights motivated by Gaussian-universality results \citep{mesforush2026lasso}, and its fits were dense, outside the sparse-signal assumption of that theory.
% src: work/datasets_v2/diagnosis/T2_T3_DIAGNOSIS.md (T2 28 of 61 wrong cells with full gold coverage, T3 13 of 52; gold-only glm-5.3-flash T2 12/18, T3 17/18)
A development diagnosis separated reader failures from retrieval failures.
With the local 4B reader, 28 of 61 wrong T2 cells had all required gold evidence in front of the reader, against 13 of 52 wrong T3 cells.
Given only the gold evidence, GLM-5.3-flash answered 12 of 18 T2 and 17 of 18 T3 development questions.
This is consistent with retrieval being a large part of the T3 failures. The development reader was therefore replaced by a stronger model.

\section{Prompts}
\label{app:prompts}
% src: RAG-Paper-upgrades/ragbench/approaches/lex5b/views.py (PROMPT); ragbench/contracts.py (ANSWER_V2_* systems)
% Blind-review fix (2026-10-07): all prompts verbatim from the frozen code tree (RAG-Paper-upgrades, commit 77924d4, clean): answer: contracts.py _V2_HEAD, SOURCE_RULE_EVIDENCE, ABSTAIN_RULE_EVIDENCE, ANSWER_V2_{SINGLE,LIST,COMPARISON,SERIES,EXPLANATION}_SYSTEM, ANSWER_USER_EVIDENCE, V2_ABSTAIN_FORMAT, LANG_NAMES; config.yaml evidence.abstain_token NOT_FOUND;
%   SetR: approaches/setr/approach.py SYSTEM, USER, render_passage; judge: tools/judge_multi.py V2_MULTI_SYSTEM, V2_MULTI_USER, V2_UNANS_SYSTEM, V2_UNANS_USER, reference_packet_v2, parse_multi, score_v2 (rubric "judge-multi v2 (strict)", MANIFEST.json judge). Python's doubled braces are shown single.
The selector, answer and judge prompts used on test are reproduced here verbatim, with placeholders in braces. The prompt of the selector of \setr{} and the judge prompt for unanswerable questions follow the other prompts.

\emph{\lexb{} selector.} In this prompt, \texttt{\{legend\}} lists the documents as \texttt{D<k>: title}, and each passage is rendered as \texttt{[n] D<k> > section}, followed by its excerpt. The parser reads ``at most 2 alternatives'' as at most two bundles per part, a first choice and one alternative (\ref{app:benchmark} of the main paper).
\begin{lstlisting}
You are selecting evidence passages for a question. Read the question and the numbered passages. Each passage is shown by its document, its section and a short excerpt; the full passage is longer.

Question: {question}

Documents:
{legend}

Passages:
{passages}

Task:
1. Split the question into the separate parts it asks for: one part per separate piece of information needed to answer it. A simple question has one part. Use at most 8 parts.
2. For each part, give the passages that together fully support its answer. Include every passage the answer depends on: each value or operand, each condition, qualifier, exception or definition - not only the passage that states the final answer. Join passages that are needed together with +. If a different set of passages would also fully support the part on its own, give it after | as an alternative (at most 2 alternatives, best first). Do not list passages that are only on a related topic. If no passage supports a part, write NONE.

Reply with exactly these lines and nothing else (numbers are passage numbers). The first line gives the number of parts, then one line per part, and the last line is the word END. An alternative must be a different set of passages.
parts=<number of parts>
P1: <passages joined by +> | <alternative passages joined by +>
P2: NONE
END
Always finish your reply with the line END.
\end{lstlisting}

\emph{Answer prompts.} All pipelines and readers use the same system prompt for a given answer format. It starts with a shared head, in which \texttt{\{answer\_lang\}} is \texttt{Persian (Farsi)} or \texttt{English}:
\begin{lstlisting}
You answer questions about documents such as company reports, books and laws.
Use only the numbered evidence passages in the user message.
If the answer is not there, reply with exactly NOT_FOUND and nothing else.
Answer in {answer_lang}.
\end{lstlisting}
The format instruction follows on the same line as the last sentence of the head. The five instructions are:
\begin{lstlisting}
[single] The question asks for one fact. Reply with the answer only, in one or two short sentences, naming what it applies to (the subject and, where the question fixes them, the period, version or condition). Copy numbers exactly as they are written, with their unit. Do not explain your reasoning.

[list] The question asks for several items. Give them as a numbered list, one short line each, naming the subject, period or version of an item whenever the question spans more than one. List each item once. Copy numbers exactly as they are written. Do not explain your reasoning.

[comparison] The question compares things across subjects, documents, periods or versions. First give each value or statement the comparison needs on its own line, in the form <subject> - <period or version>: <value or statement>, and write NOT_FOUND for one you cannot give. Then give the comparison, difference or calculation the question asks for in one or two short sentences, with any arithmetic in brief. If the definitions, periods or units differ, say that the values are not directly comparable instead of comparing them. Copy numbers exactly as they are written. Do not explain your reasoning.

[series] The question asks for one value over several periods or versions. Reply with one line per period or version, earliest first, in the form <period or version>: <value> <unit>, and write <period or version>: NOT_FOUND for one whose value you cannot give. Copy numbers exactly as they are written. End with one short sentence on how the value changed; if what is counted differs between periods (a different definition or scope), say instead that the values are not directly comparable. Do not explain your reasoning.

[explanation] The question needs an answer that combines several points (for example a rule and its exception or condition, a definition applied to a case, a cause and its effect, or several steps of one account). Give the direct answer first, then each point it rests on, in at most six short sentences, naming the article, section, subject or period each point comes from. Answer every part of the question. Copy numbers exactly as they are written. Do not add points you cannot support.
\end{lstlisting}
The bracketed labels were added for this paper and are not sent. The user message is \texttt{Evidence:} followed by the rendered pack, a blank line and \texttt{Question: \{question\}}. Unanswerable questions are given the format of the tier they imitate (T1 single, T2 explanation, T3 comparison) and are correct only with an abstention. Completion limits are given in Table~\ref{tab:settings-full}.
\emph{\setr{} selector.} The prompt is the system message and the selection prompt of the SetR release (\texttt{generate\_data.py}, selection mode with information requirements), verbatim. Only the number of passages, the question and the passages are filled in. Each of the top 20 passages of \hybrid{} is rendered as \texttt{[i]} followed by its text, with line breaks replaced by spaces and the text cut at 1,500 characters, and the passages are joined by line breaks. The system message is:
\begin{lstlisting}
You are RankLLM, an intelligent assistant that can rank passages based on their relevancy to the query.
\end{lstlisting}
The user message is:
\begin{lstlisting}
I will provide you with {num} passages, each indicated by a numerical identifier []. Select the passages based on their relevance to the search query: {question}.

{context}

Search Query: {question}

Please follow the steps below:
Step 1. Please list up the information requirements to answer the query.
Step 2. for each requirement in Step 1, find the passages that has the information of the requirement.
Step 3. Choose the passages that mostly covers clear and diverse informations to answer the query. Number of passages is unlimited. The format of final output should be '### Final Selection: [] []', e.g., ### Final Selection: [4] [2].
\end{lstlisting}
The parser takes the numbers after the last \texttt{Final Selection} marker, keeps the first occurrence of each valid number and keeps the order of the model. A reply without a marker or a valid number, a missed deadline of 8 seconds and a call error all fall back to \hybrid{} order, with no repair and no retry. The selected passages are packed first, and the other candidates of \hybrid{} fill the rest of the budget.

\emph{Judge.} The judge grades every answer of an answerable question with the following system and user prompts:
\begin{lstlisting}
[system] You are a strict, impartial grader for a question-answering benchmark over documents (company reports, non-fiction books and laws) written in Persian and English. You compare a candidate answer with a reference packet built from the documents: the reference answer, the answer requirements (each with an id, marked required or optional), the source excerpts behind them and accepted alternatives. The packet is the truth; do not use outside knowledge. Reply with one JSON object only.

[user] Question ({lang}): {question}
Answer format: {fmt}

Reference packet:
{packet}

Candidate answer:
{answer}

Grade in two lists.
1. "claims": split the candidate into its distinct claims. For each give {"text": <the claim, short>, "unit":
   <the requirement id it expresses, or null>, "status": one of correct | wrong_value | wrong_binding | contradiction |
   unsupported | extra | duplicate | non_substantive}.
   correct: states a requirement with the same meaning, bound to the same subject, aspect and constraints (time,
   version, place, population, condition). Paraphrase, either language, rounding at the requirement's precision,
   Persian digits, the Persian decimal slash (78/3 = 78.3) and scale words are all fine.
   wrong_value: the right subject and aspect with a different value or statement. wrong_binding: a right value or
   statement attached to the wrong subject, time, version or condition. contradiction: conflicts with a requirement,
   an excerpt or another claim of the answer. unsupported: a substantive claim not in the packet that changes the
   answer (a different conclusion, a wrong reason, rule, exception or trend). extra: harmless additional detail not in
   the packet that does not change the answer. duplicate: repeats an earlier claim. non_substantive: framing, hedging or
   citation markers such as [2].
2. "units": for EVERY requirement id in the packet give one of correct | missing | wrong | contradicted.
A derived requirement (a comparison, difference, ranking or conclusion) is correct only if the candidate states it;
its inputs are graded as their own requirements. For a list, each required item is its own requirement.
Reply with JSON only: {"claims": [...], "units": {"<requirement id>": "<status>", ...}, "reason": "at most 25 words"}
\end{lstlisting}
The reference packet is a JSON object built from the question file alone, with the question's tier, answer format, integration operation, language and reference answer, and one entry per requirement (identifier, whether required, proposition, subject, aspect, constraints, value, unit, precision, accepted alternatives, derivation and source excerpts). A reply is valid only if its \texttt{units} grade every requirement identifier of the packet and its \texttt{claims} use the eight claim statuses; an invalid reply is retried up to three attempts and never repaired. An answer is strictly correct if every required requirement is graded \texttt{correct}, none is \texttt{contradicted}, and no claim is \texttt{contradiction}, \texttt{wrong\_value}, \texttt{wrong\_binding} or \texttt{unsupported}. Claims marked \texttt{extra}, \texttt{duplicate} or \texttt{non\_substantive}, and optional requirements, do not affect it.

For unanswerable questions the judge receives the absence audit of the question instead of a packet. An answer that is exactly the abstention token \texttt{NOT\_FOUND} is ruled correct without a call. Any other answer is graded with the following prompts, in which \texttt{\{missing\}} is what the documents lack, and \texttt{\{near\}} lists the nearby facts that are not the answer, joined by semicolons:
\begin{lstlisting}
[system] You are a strict, impartial grader for a question-answering benchmark over documents (company reports, non-fiction books and laws) written in Persian and English. The question asks for something the documents do not contain. Reply with one JSON object only.

[user] Question ({lang}): {question}
What the documents lack: {missing}
Nearby facts that are NOT the answer: {near}

Candidate answer:
{answer}

The candidate is CORRECT only if it declines to give the missing fact (for example NOT_FOUND, "the documents do not
say", or an explanation that the value is not given), possibly while mentioning the nearby facts as such. It is
INCORRECT if it states a value, date, name or conclusion for the missing fact, including one taken from a nearby fact.
Reply with JSON only: {"correct": true or false, "reason": "at most 20 words"}
\end{lstlisting}
The reply is a JSON object holding a verdict and a reason of at most 20 words.

\section{Examination of the gold-evidence misses}
\label{supp:goldmiss}
% src: work/posthoc_2026-11/S17_goldmiss/RESULTS.md (the table, the counts and the reasons); S2a_packs/RESULTS.md sec. 2.4 (83 of 90; T1 27/30, T2 30/30, T3 26/30; en 40/45, fa 43/45); S8_factorial/tables.md 'Accuracy, judge 2' (GOLD 82..84, 88 resolved)
With the first annotated bundle alone as evidence, the primary reader answered 83 of the 90 answerable test questions correctly under the first judge. Table~\ref{tab:goldmiss} lists the seven misses, with the tier, the language, the expected answer, the given answer, the verdicts of both judges and a class. The pack that the reader saw for each question is the quoted bundle in bundle order, and each class is one reading of the saved question, pack, answer and judge records. No second person checked it.

Four of the seven packs do not contain enough text to tie the key value to its subject. In three of them the quote is a table row or a clause that is cut out of its context, and in the fourth the passage that fixes the period is missing. In these cases the reader reported that it could not give the value, as its prompt asks. Two misses are judge decisions that a person would probably not make. One was failed by the first judge on a unit that the source itself gives, the other because the first judge did not link a correct line to its requirement. The second judge accepts the second answer and fails the first on a restated source sentence. One miss is an answer that is true but leaves out one of two required points, and the saved data do not show whether a person would accept it.
The second judge scored 88 of the 90 answers and found 82 correct. It leaves two of the seven unresolved because its reply used a claim status that the frozen parser rejects, and it fails two answers that the first judge accepted. If the two judge cases were counted as correct, the gold reference would be 85 of 90. This is a count from the reading above and not a new score.
% src: work/posthoc_2026-11/S17_goldmiss/RESULTS.md (one row per question; classes are one reading of the saved records, no second annotator).
%      Reader answers: S2a_packs/answers.jsonl (arm GOLD) and S2a_packs/packs/GOLD/<id>.json; judge 1: S2a_packs/judge/GOLD/verdicts_openai__gpt-oss-120b.jsonl (83 of 90);
%      judge 2: S8_factorial/judge2/GOLD/verdicts_mistralai__mistral-large-2512.jsonl (88 resolved, 82 correct; RPT-en-T3-002 and LAW-en-T1-004 unresolved);
%      questions and gold bundles: datasets_v2/test_freeze/questions/test_v2run_all.jsonl. Counts: S2a_packs/RESULTS.md sec. 2.4 (T1 27/30, T2 30/30, T3 26/30; en 40/45, fa 43/45).
\begin{table*}[tp]
\caption{Post hoc. The seven answerable test questions that the reader missed with the gold bundle. Persian answers are given in English}
\label{tab:goldmiss}
\centering
\scriptsize
\setlength{\tabcolsep}{2.5pt}
\begin{tabularx}{\textwidth}{@{}>{\raggedright\arraybackslash}p{1.75cm}cc>{\raggedright\arraybackslash}X>{\raggedright\arraybackslash}X>{\raggedright\arraybackslash}p{1.15cm}>{\raggedright\arraybackslash}p{1.15cm}>{\raggedright\arraybackslash}p{4.6cm}@{}}
\toprule
Question & Tier & Lang. & Expected answer & Given answer & Judge 1 & Judge 2 & Class and reason \\
\midrule
RPT-en-T3-002 & T3 & en & ArcelorMittal 57.9 Mt in 2024, Tata Steel Group 29.94 Mt in FY2023-24, ratio about 1.9 & Tata value correct; ArcelorMittal ``NOT\_FOUND''; no ratio & wrong & unresolved & \emph{Bundle not sufficient.} The pack holds 57.9 only as the bare row ``Total 43.5 14.4 57.9'' and a separate caption, so nothing says it is the company total. \\
RPT-en-T3-001 & T3 & en & Mowi 501,530 t and Ler\o y 147,701 t of salmon, ratio about 3.4 & Mowi value correct; Ler\o y ``NOT\_FOUND'' & wrong & wrong & \emph{Bundle not sufficient.} The Ler\o y figure is the row ``Salmon Trout $|$ 147 701'' with its column header in another item, and the total row that would resolve it is not in the bundle. \\
BOOK-en-T1-005 & T1 & en & Population starts at about 7 million; after fifty years 28 million, subsistence for 21 million & Starts at 7 million; gives 28 and 21 million for ``after twenty-five years''; fifty years ``NOT\_FOUND'' & wrong & wrong & \emph{Bundle not sufficient.} The quoted sentence says ``in the next twenty-five years''. The passage that makes it the second period is not in the bundle. \\
BOOK-fa-T3-002 & T3 & fa & Vosough al-Dowleh was first vice-president of the first Majlis, concluded the 1919 agreement, and 130,000 lira was divided among three men & The same three facts, each with a source note & wrong & wrong & \emph{Judge strictness.} Every requirement is present. Judge 1 faulted the unit, which the source gives as lira. Judge 2 graded all requirements correct and failed a restated source sentence as unsupported. \\
LAW-fa-T1-002 & T1 & fa & Five minority seats: Zoroastrians 1, Jews 1, Assyrian and Chaldean Christians 1, Armenian Christians 1 each in the south and the north & Total of 5 with the same five lines & wrong & correct & \emph{Judge strictness.} Judge 1 marked the Jews requirement missing, although the answer states it, and gave the line no requirement. Judge 2 accepts the answer. \\
LAW-en-T1-004 & T1 & en & ``Without undue delay'' and in any event within 30 days (Consumer Rights Act 2015, s.28(3)) & Only the 30-day limit & wrong & unresolved & \emph{Reader error.} The pack holds both limbs. The answer is true but omits the first required point. \\
LAW-en-T3-004 & T3 & en & Police force: higher maximum amount, \pounds 17,500,000; EU company: EUR 20 million or 4\% of turnover & Higher maximum amount ``not stated''; EU figure correct & wrong & wrong & \emph{Bundle not sufficient.} The \pounds 17.5 million sentence appears before its parent line ``The higher maximum amount is'', and the pack lacks the clause for an undertaking. \\
\bottomrule
\end{tabularx}
\end{table*}


\FloatBarrier
\bibliographystyle{elsarticle-num}
\bibliography{refs}
\end{document}
